\section{Trayectoria electrónica, espectro de retorno y conservación de la masa}
\label{vi:sec:trayectoria-espectral-masa}

Esta sección desarrolla, sobre el operador de masa ya construido, la parte
dinámica que conserva la trayectoria electrónica. No introduce una segunda
definición del electrón ni una calibración adicional. La cadena causal es la
misma del artículo:
\[
 \mathrm{APP}\longrightarrow\mathrm{TRIT}\longrightarrow\mathrm{TPK}
 \longrightarrow X_{\rm enr}
 \longrightarrow\mathcal C_{\rm cont}^{\rm disc}
 \longrightarrow (\Gamma_e,\tau_e,\sigma_e)
 \longrightarrow M^\beta .
\]
Las coordenadas $\pi_{\rm HMT},\varphi_{\rm HMT},e_{\rm HMT}$ y
$\alpha_{\rm HMT}$ son salidas HMT correlacionadas generadas desde el estado
enriquecido; no se introducen como datos convencionales. El análisis de Fourier
que aparece más abajo se aplica a la palabra ya producida y no selecciona la ruta.

\subsection{Objetos centrales distintos y memoria de la fibra}

La celda fija $(5,5)$, el dato de frontera $111111$, la emisión transversal
$555555$ y la orientación $\tau_e=(+++---)^2$ no son nombres de un mismo
objeto. La involución celular
\[
 \rho_c(r,c)=(10-r,10-c)
\]
tiene un único punto fijo, $(5,5)$. En él,
\[
 5+5=1+9,\qquad 5\cdot5=7+2\cdot9,
\]
mientras que en $(8,2)$ y $(2,8)$ se conserva la suma pero cambia el cociente
del producto. La celda fija determina la excepción central; el registro de
frontera y la emisión transversal conservan otras coordenadas del estado.

En particular, la fibra de $555555$ en el canal de producto contiene
$24$ posiciones orientadas. Tras un avance del TPK se distribuye en
\[
 555177,\qquad 555717,\qquad 555771,
\]
con ocho posiciones por salida. El siguiente valor no es función de la firma
escalar aislada: para reconstruir la evolución deben conservarse la clase de
memoria y la orientación. Esta distinción es la razón por la que el operador
de masa actúa sobre fibras de rutas y no sobre una lista de cifras.

\subsection{Soporte celular exacto de la trayectoria central}

La trayectoria de~\eqref{vi:dep:trayectoria-electron} comienza en
$(r,c,h,v)=(5,5,1,1)$. Cada bloque TRIT $(+1,-1,0)$ avanza una vez cada
coordenada y conserva un paso de umbral. En los extremos de cada bloque se
tiene $r=c$; en el paso intermedio,
\[
 c-r\equiv h\pmod 9.
\]
Nueve bloques recorren todas las posiciones de la diagonal correspondiente.
Las inversiones simultáneas de $(h,v)$ en los pasos $27,54,81$ añaden la
diagonal conjugada sin reiniciar el estado. Por consiguiente, la proyección
celular de los $108$ pasos tiene el soporte exacto
\begin{equation}
 \operatorname{Supp}_{\rm cell}(\Gamma_e)
 =\left\{(r,c)\in\{1,\ldots,9\}^2:
 c-r\equiv0,1,-1\pmod9\right\},
 \qquad
 \left|\operatorname{Supp}_{\rm cell}(\Gamma_e)\right|=27.
 \label{vi:eq:soporte-electron-27}
\end{equation}
La igualdad se obtiene por la regla de actualización, no por un conteo
retrospectivo. Describe una morfología discreta de trayectoria en el toro
posicional. Su realización como banda de Möbius requiere además el transporte
de orientación y el pegado de la sección~\ref{vi:sec:conjugacion}; no se
identifica el soporte finito con una forma material tridimensional.

\subsection{Dos historias, una orientación sectorial y reconstrucción exacta}

La biografía central contiene las historias
\begin{equation}
 e_{++}=(2,2,5,-4,-3,-2)^2,
 \qquad
 e_{--}=(4,3,2,-2,-2,-5)^2.
 \label{vi:eq:historias-electronicas}
\end{equation}
Ambas se expresan en la coordenatización sectorial fija
$\Sigma_{12}=(+++---)^2$. La orientación absoluta de la recurrencia se
transporta por separado; no se identifica el signo de una historia con la
etiqueta sectorial. Las dos historias tienen suma cero y el mismo histograma
de valores absolutos, pero no la misma memoria posicional.

Para una historia $e=(e_0,\ldots,e_{11})$, definimos
\begin{equation}
 p_i=e_i+e_{i+6},\qquad a_i=e_i-e_{i+6},
 \qquad s_C=\sum_{i=0}^{5}p_i,
 \qquad M_i=p_i-p_5\quad(0\leq i<5).
 \label{vi:eq:codificador-memoria-electron}
\end{equation}
El decodificador es
\begin{equation}
 p_5=\frac{s_C-\sum_{i=0}^{4}M_i}{6},\quad
 p_i=p_5+M_i,\quad
 e_i=\frac{p_i+a_i}{2},\quad e_{i+6}=\frac{p_i-a_i}{2}.
 \label{vi:eq:decodificador-memoria-electron}
\end{equation}
Las condiciones de integridad son la divisibilidad por seis de la primera
expresión y la paridad común de $p_i$ y $a_i$. Para
\eqref{vi:eq:historias-electronicas} se obtiene
\[
 M^{++}=(8,8,14,-4,-2),\qquad
 M^{--}=(18,16,14,6,6),\qquad a=0.
\]
Sustituir estos datos en~\eqref{vi:eq:decodificador-memoria-electron}
restituye las doce componentes. Queda así probada la distinción de historias
aunque una lectura parcial coincida. No se concluye que sus masas deban ser
distintas: un lector de masa puede ser invariante bajo un transporte que
conserva menos coordenadas que el estado enriquecido.

\subsection{Identidad espectral del lector de retorno}

Sea $\tau=(\tau_0,\ldots,\tau_{11})\in\{-1,+1\}^{12}$ y sean los índices
cíclicos. Definimos
\begin{equation}
 C_s(\tau)=\sum_{j=0}^{11}\tau_j\tau_{j+s},
 \qquad
 \Omega(\tau)=-2C_1-4C_2-6C_3,
 \qquad
 P_6(\tau)=\frac{C_6}{2}.
 \label{vi:eq:lectores-correlacion}
\end{equation}
La primera igualdad es la forma desarrollada de
$\Omega=4A_3-3A_4$ ya demostrada en
\eqref{vi:dep:KOmega}. Para hacer visible qué modos transportan la
corrección, introducimos, después de producir $\tau$,
\begin{equation}
 T_k(\tau)=\sum_{j=0}^{11}\tau_j
 e^{-2\pi i k j/12},\qquad 0\leq k<12.
 \label{vi:eq:dft-retorno}
\end{equation}

\begin{proposition}[Descomposición espectral exacta]
Para toda palabra real de longitud doce,
\begin{align}
 C_s&=\frac1{12}\sum_{k=0}^{11}|T_k|^2
          \cos\!\left(\frac{2\pi ks}{12}\right),\label{vi:eq:wiener-khinchin-finito}\\
 \Omega&=\frac1{12}\sum_{k=0}^{11}|T_k|^2w_k,
 \quad
 w_k=-2\cos\theta_k-4\cos2\theta_k-6\cos3\theta_k,
 \quad\theta_k=\frac{2\pi k}{12},\label{vi:eq:omega-espectral}\\
 P_6&=\frac1{24}\sum_{k=0}^{11}|T_k|^2(-1)^k.
 \label{vi:eq:p6-espectral}
\end{align}
\end{proposition}
\begin{proof}
Al desarrollar $|T_k|^2$ aparece
$\sum_{j,\ell}\tau_j\tau_\ell e^{-2\pi ik(j-\ell)/12}$.
La suma sobre $k$ vale doce cuando $j-\ell$ coincide con el desplazamiento
prescrito y cero en otro caso. Ésta es la ortogonalidad finita de caracteres
y demuestra~\eqref{vi:eq:wiener-khinchin-finito}. Insertar sucesivamente
$s=1,2,3$ en~\eqref{vi:eq:lectores-correlacion} produce
\eqref{vi:eq:omega-espectral}; $s=6$ produce
\eqref{vi:eq:p6-espectral}.
\end{proof}

Para $\tau_e=(+++---)^2$ puede escribirse, con
$z_k=e^{-2\pi ik/12}$,
\begin{equation}
 T_k=(1+(-1)^k)(1-z_k^3)(1+z_k+z_k^2).
 \label{vi:eq:factorizacion-tk-electron}
\end{equation}
Sólo sobreviven $k=2,6,10$, con
\[
 T_2=4-4i\sqrt3,\qquad T_6=4,\qquad
 T_{10}=4+4i\sqrt3,
\]
y, por tanto,
\begin{equation}
 |T_2|^2=64,\qquad |T_6|^2=16,\qquad |T_{10}|^2=64.
 \label{vi:eq:potencias-espectrales-electron}
\end{equation}
Los pesos $w_2,w_6,w_{10}$ son $7,4,7$. En consecuencia,
\begin{equation}
 \Omega(\tau_e)=\frac{64\cdot7+16\cdot4+64\cdot7}{12}=80,
 \qquad P_6(\tau_e)=6.
 \label{vi:eq:omega-p6-electron}
\end{equation}
Esto recupera el segundo lector de~\eqref{vi:dep:KOmega} desde los modos
de la trayectoria. La palabra alternante $(+-)^6$, que conserva el mismo
histograma de signos, tiene $|T_6|^2=144$ y produce $\Omega=48$; constituye
un control negativo de que el orden, y no sólo el conteo, interviene.
Las traslaciones cíclicas multiplican $T_k$ por una fase y conservan
$|T_k|^2$, de modo que el lector no recupera por sí solo el origen temporal.

La contribución de retorno al exponente electrónico queda expresada por
\begin{equation}
 \kappa_e=
 \frac1{12}\sum_{k=0}^{11}|T_k|^2
 \left(w_k+\frac{\Delta_4}{2}(-1)^k\right)-54\alpha_{\rm HMT}
 =80-54\alpha_{\rm HMT}+6\Delta_4.
 \label{vi:eq:kappa-espectral-electron}
\end{equation}
Esta identidad no introduce un Hamiltoniano espacial: demuestra que el
corrector de masa ya definido es un funcional espectral exacto de la palabra
de retorno.

\subsection{Extensión de la ecuación electrónica a las multisecciones}

El registro de una ruta $\gamma$ suministra
$\sigma_\gamma=(n_A,s_C,k_\Delta,b_{120},b_\zeta)$ y las coordenadas de
torre $\eta_\gamma=(\eta_h)_h$. En una coordenatización común, el carácter refinado es
\begin{equation}
 \begin{aligned}
 \log\chi_{\rm ref}(\sigma_\gamma,\eta_\gamma)
 &=n_Ax+\frac{s_C}{6}y+k_\Delta\Delta_4+b_\zeta\zeta_{270}
   +b_{120}s_{120}+\sum_{h\in\langle90,120\rangle}\eta_hs_h\\
 &=n_Ax+\frac{s_C}{6}y+k_\Delta\Delta_4+b_\zeta\zeta_{270}
   +\sum_{h\in\langle90,120\rangle}N_hs_h,\\
 N_h&:=\eta_h+b_{120}\mathbf1_{\{h=120\}}.
 \end{aligned}
 \label{vi:eq:caracter-refinado-explicito}
\end{equation}
donde $x=A_{\rm rad}$, $y=C^*_{\rm rad}$,
$\zeta_{270}=135\Delta_4^2$ y $s_h=q_-^h-q_+^h$. Las dos líneas son la
misma evaluación: $b_{120}$ aparece una sola vez, como término explícito o
absorbido en $N_{120}$. Para una torre infinita se exige convergencia absoluta
de la última serie. Esta fórmula conserva separadas la torsión $\zeta_{270}$,
las dos hojas y las dos composiciones de altura $360=4\cdot90=3\cdot120$.

En la base de rutas de una fibra, sea
$B^r e_\gamma=\kappa_\gamma^r e_\gamma$, con $r=D$ u $\Omega$. La ley
operatoria de la sección~\ref{vi:sec:operador-masa} da, en su coordenatización diagonal,
\begin{equation}
 \boxed{
 \frac{m_\gamma^\beta}{m_e^\beta}=
 \chi_{\rm ref}(\sigma_\gamma-\sigma_e,
                 \eta_\gamma-\eta_e)
 R_{\rm act}^{\kappa_\gamma-\kappa_e}.}
 \label{vi:eq:ley-masas-trayectoria}
\end{equation}
La ecuación electrónica es la sección de referencia de esta identidad, no
una cifra experimental usada para escoger rutas. Las etiquetas físicas se
asignan después de construir fibras, representaciones y canales; la semejanza
numérica no sustituye ese emparejamiento prospectivo.

\subsection{Criterio operatorio de transporte a masa fija}

Definimos la coordenada adimensional
\begin{equation}
 \Lambda_\gamma:=\log\frac{m_\gamma^\beta}{m_e^\beta}
 \label{vi:eq:lambda-masa-ruta}
\end{equation}
en una coordenatización común con masas positivas. Este símbolo no es la longitud
recíproca $\Lambda=L_\alpha X^{-1}$ ni la coordenada
$\bar\lambda_\Gamma$ empleada en otras secciones.

Sea $\mathscr D$ un dominio finito de rutas estable bajo un transporte TPK
biyectivo $F:\mathscr D\to\mathscr D$, y sea
$V_Fe_\gamma=e_{F\gamma}$ su operador de permutación. Entonces
\begin{equation}
 [M^\beta,V_F]e_\gamma
 =m_e^\beta\left(e^{\Lambda_{F\gamma}}-e^{\Lambda_\gamma}\right)
 e_{F\gamma}.
 \label{vi:eq:conmutador-masa-transporte}
\end{equation}
Por tanto,
\begin{equation}
 [M^\beta,V_F]=0
 \quad\Longleftrightarrow\quad
 \Lambda_{F\gamma}=\Lambda_\gamma
 \quad\text{para toda }\gamma\in\mathscr D.
 \label{vi:eq:criterio-conservacion-masa}
\end{equation}
Usando~\eqref{vi:eq:caracter-refinado-explicito}, la condición es
\begin{equation}
 \delta n_Ax+\frac{\delta s_C}{6}y
 +\delta k_\Delta\Delta_4+\delta b_\zeta\zeta_{270}
 +\sum_h\delta N_hs_h
 +\delta\kappa\log R_{\rm act}=0.
 \label{vi:eq:balance-conservacion-masa}
\end{equation}
Aquí
\(
 \delta N_h=\delta\eta_h+
 \delta b_{120}\mathbf1_{\{h=120\}},
\)
de modo que la diferencia conserva las coordenadas de torre originales y
añade el término de incidencia de $b_{120}$ exactamente una vez.
Se comprueba después de producir $F$ mediante el TPK; no se selecciona
$F$ retrospectivamente por el valor de la masa. El criterio permite que
cambien eventos y memoria siempre que se conserve la coordenada total.
No afirma que todo transporte TPK conserve la masa ni que agote la dinámica
de un sistema ligado. Sus hipótesis son: dominio estable, transporte
biyectivo, coordenatización fija, masas positivas y convergencia de las torres.

\subsection{Resultado conjunto}

Las ecuaciones~\eqref{vi:eq:soporte-electron-27},
\eqref{vi:eq:decodificador-memoria-electron},
\eqref{vi:eq:kappa-espectral-electron},
\eqref{vi:eq:ley-masas-trayectoria} y
\eqref{vi:eq:criterio-conservacion-masa} cierran una misma cadena. La
trayectoria central recorre exactamente veintisiete celdas; sus historias
conservan memoria reversible; el retorno posee tres modos no nulos que
producen $(\Omega,P_6)=(80,6)$; la ley de masas extiende la sección
electrónica a las multisecciones; y el conmutador identifica exactamente
qué transportes preservan la masa. Morfología celular, espectro de retorno y
magnitud son lectores coordinados del mismo estado enriquecido, no
observables intercambiables.
